%==============================================================================
\appendix
\section{Rendezvous Agent --- GLP Implementation}
\label{appendix:rv-agent}
%==============================================================================
The rendezvous agent processes cold-call messages on its input stream.  A $\mathit{reconnect}$ message from agent~$A$ includes~$A$'s public key, friend~$B$'s public key, a friendship proof, and a reply variable.  An $\mathit{available}$ message from~$B$ includes both public keys and a reply variable.

The agent verifies the friendship proof via \code{verify\_friendship} and observes each connecting agent's public address via \code{peer\_address} (both body system predicates).  It searches a pending list for a matching request from the other side.  If found, \code{complete\_rendezvous} constructs the reply values in its head: $A$ receives $\mathit{punch}(\mathit{Addr}_B)$ and $B$ receives $\mathit{punch}(\mathit{Addr}_A)$.  If the proof is invalid, the reply is bound to $\mathit{rejected}$.  If no match is found, the request is stored in the pending list.

The key GLP mechanism is the \emph{reply-variable chain}: each cold-call message contains a reply variable (a writer held by the sender).  As the rendezvous agent processes the message, the reply variable is threaded through writer$\to$reader handoffs.  When both sides are matched, the assignments propagate back through the chain to the original senders.

\begin{verbatim}
SignalReply ::= punch(Constant) ; rejected.
VerifyResult ::= ok ; invalid.

SignalMsg ::= reconnect(Constant, Constant, Constant, SignalReply?)
              %% reconnect(A_pk, B_pk, Proof, Reply)
            ; available(Constant, Constant, SignalReply?).
              %% available(B_pk, A_pk, Reply)

PendingEntry ::= needs(Constant, Constant, Constant, SignalReply?)
                 %% needs(A_pk, B_pk, AddrA, ReplyA?)
               ; ready(Constant, Constant, Constant, SignalReply?).
                 %% ready(B_pk, A_pk, AddrB, ReplyB?)
PendingList ::= [] ; [PendingEntry | PendingList].

procedure rv_agent(Stream(SignalMsg)?, PendingList?).

rv_agent([reconnect(A, B, Proof, ReplyA?)|In], Pending) :-
    ground(A?), ground(B?) |
    verify_friendship(Proof?, A?, B?, Status),
    peer_address(A?, AddrA),
    handle_verify(Status?, A?, B?, AddrA?, ReplyA, Pending?, Pending1),
    rv_agent(In?, Pending1?).

rv_agent([available(B, A, ReplyB?)|In], Pending) :-
    ground(A?), ground(B?) |
    peer_address(B?, AddrB),
    find_needs(A?, B?, AddrB?, ReplyB, Pending?, Pending1),
    rv_agent(In?, Pending1?).

rv_agent([], _).

procedure handle_verify(VerifyResult?, Constant?, Constant?, Constant?,
                        SignalReply, PendingList?, PendingList).
handle_verify(ok, A, B, AddrA, ReplyA?, Pending, Pending1?) :-
    find_ready(A?, B?, AddrA?, ReplyA, Pending?, Pending1).
handle_verify(invalid, _, _, _, rejected, Pending, Pending?).

procedure find_ready(Constant?, Constant?, Constant?, SignalReply,
                     PendingList?, PendingList).
find_ready(A, B, AddrA, ReplyA?, [ready(B2, A2, AddrB, ReplyB?)|Ps], Ps?) :-
    A? =?= A2?, B? =?= B2? |
    complete_rendezvous(AddrA?, AddrB?, ReplyA, ReplyB).
find_ready(A, B, AddrA, ReplyA?, [P|Ps], [P?|Ps1?]) :-
    otherwise |
    find_ready(A?, B?, AddrA?, ReplyA, Ps?, Ps1).
find_ready(A, B, AddrA, ReplyA?, [], [needs(A?, B?, AddrA?, ReplyA)]).

procedure find_needs(Constant?, Constant?, Constant?, SignalReply,
                     PendingList?, PendingList).
find_needs(A, B, AddrB, ReplyB?, [needs(A2, B2, AddrA, ReplyA?)|Ps], Ps?) :-
    A? =?= A2?, B? =?= B2? |
    complete_rendezvous(AddrA?, AddrB?, ReplyA, ReplyB).
find_needs(A, B, AddrB, ReplyB?, [P|Ps], [P?|Ps1?]) :-
    otherwise |
    find_needs(A?, B?, AddrB?, ReplyB, Ps?, Ps1).
find_needs(A, B, AddrB, ReplyB?, [], [ready(B?, A?, AddrB?, ReplyB)]).

procedure complete_rendezvous(Constant?, Constant?, SignalReply, SignalReply).
complete_rendezvous(AddrA, AddrB, punch(AddrB?), punch(AddrA?)).
\end{verbatim}

The system predicates \code{peer\_address/2}, \code{verify\_friendship/4}, and \code{punch\_udp/1} are described in Section~\ref{sec:rv-predicates}.

\paragraph{Open: cold-call branch not yet supported.}
The code above handles only the friend-reconnect form of \code{reconnect}: every \code{reconnect} message is fed unconditionally to \code{verify\_friendship/4}, which rejects anything that is not a valid friendship attestation.  The IP cold-call flow (Appendix~\ref{appendix:ip-cold-call-invites}, step~5), however, has~$A$ send a \code{reconnect} whose \code{Proof} slot carries an invite-ID reference to the pre-registered invite rather than a friendship attestation; the rendezvous agent is expected to dispatch on the proof variant and run a different verifier.  Two candidate resolutions, neither applied here:
\begin{itemize}
\item Make the \code{Proof} type polymorphic (friendship attestation vs.\ invite-ID reference) and have \code{rv\_agent} dispatch on the variant before calling the appropriate verifier.
\item Introduce a dedicated \code{cold\_call\_accept} message type with its own handler, leaving \code{reconnect} untouched.
\end{itemize}
The code in this appendix is left as-is pending that choice.

%==============================================================================
\section{Rendezvous Client --- GLP Implementation}
\label{appendix:rv-client}
%==============================================================================
The client side provides building blocks for agents that need to use a rendezvous agent.  \code{make\_reconnect} constructs a $\mathit{reconnect}$ cold-call message with a fresh reply variable embedded, returning the message and the paired reply reader.  \code{make\_available} does the same for $\mathit{available}$ messages.  \code{handle\_reply} dispatches on the reply: $\mathit{punch}(\mathit{Addr})$ triggers \code{punch\_udp}; $\mathit{rejected}$ requires no networking action.

\begin{verbatim}
SignalReply ::= punch(Constant) ; rejected.

SignalMsg ::= reconnect(Constant, Constant, Constant, SignalReply?)
              %% reconnect(A_pk, B_pk, Proof, Reply)
            ; available(Constant, Constant, SignalReply?).
              %% available(B_pk, A_pk, Reply)

procedure make_reconnect(Constant?, Constant?, Constant?,
                         SignalMsg, SignalReply).
make_reconnect(A, B, Proof,
               reconnect(A?, B?, Proof?, Reply),
               Reply?).

procedure make_available(Constant?, Constant?, SignalMsg, SignalReply).
make_available(B, A,
               available(B?, A?, Reply),
               Reply?).

procedure handle_reply(SignalReply?).
handle_reply(punch(Addr)) :- ground(Addr?) | punch_udp(Addr?).
handle_reply(rejected).
\end{verbatim}

An end-to-end test wires the client and server together on a shared message stream:
\begin{verbatim}
procedure test_full_scenario(SignalReply, SignalReply).
test_full_scenario(ReplyA?, ReplyB?) :-
    make_reconnect(alice, bob, proof, MsgA, ReplyA),
    make_available(bob, alice, MsgB, ReplyB),
    rv_agent([MsgA?, MsgB?], []).
\end{verbatim}
After execution, $\mathit{ReplyA}$ is bound to $\mathit{punch}(\mathit{stub\_addr\_bob})$ and $\mathit{ReplyB}$ to $\mathit{punch}(\mathit{stub\_addr\_alice})$.

%==============================================================================
\section{Friends-Based Rendezvous --- GLP Implementation}
\label{appendix:fmrv}
%==============================================================================
Friends-based rendezvous is implemented as an extension of the social graph agent.  The friends list is extended to include each friend's current IP address: entries have the form $\mathit{fe}(\mathit{Name}, \mathit{OutputStream}, \mathit{Addr})$.  Agents notify friends of address changes via $\mathit{addr\_update}$ messages.

When agent~$A$ asks friend~$B$ to mediate reconnection to~$C$, the procedure \code{lookup\_and\_send} performs a fused lookup-and-send: in a single traversal of the friends list, it locates a friend's entry, prepends a $\mathit{punch\_to}$ message to that friend's output stream, and returns the friend's address.  The $\mathit{punch\_to}$ argument carries an $\mathit{AddrResult}$: $\mathit{found}(\mathit{Addr})$ if the peer was located, or $\mathit{not\_found}$ if not.  The recipient dispatches accordingly --- \code{punch\_udp} is never called with an invalid address.

The $\mathit{mediate}$ clause calls \code{lookup\_and\_send} twice with cross-references: the first call sends $A$ a $\mathit{punch\_to}(\mathit{AddrC}?)$ where $\mathit{AddrC}$ is not yet bound; the second call binds $\mathit{AddrC}$ when it finds~$C$'s entry.  GLP's suspension mechanism fills in the unbound reader once the writer is bound --- the standard futures/promises mechanism.

\begin{verbatim}
AddrResult ::= found(Constant) ; not_found.

MsgOp ::= mediate(Constant)        %% mediate(C_pk)
        ; punch_to(AddrResult)     %% punch_to(found(Addr)) or punch_to(not_found)
        ; addr_update(Constant).   %% addr_update(NewAddr)

Msg ::= msg(Constant, Constant, MsgOp).
MsgStream ::= [] ; [Msg | MsgStream].
FriendStream ::= [] ; [Msg | FriendStream].
FriendEntry ::= fe(Constant, FriendStream?, Constant).
FriendList ::= [] ; [FriendEntry | FriendList].

procedure fmrv(Constant?, MsgStream?, FriendList?).

fmrv(Id, [msg(From, To, addr_update(Addr))|In], Fs) :-
    ground(Id?), Id? =?= To? |
    update_addr(From?, Addr?, Fs?, Fs1),
    fmrv(Id?, In?, Fs1?).

fmrv(Id, [msg(From, To, mediate(C))|In], Fs) :-
    ground(Id?), ground(From?), ground(C?), Id? =?= To? |
    lookup_and_send(From?, Id?, AddrC?, AddrFrom, Fs?, Fs1),
    lookup_and_send(C?, Id?, AddrFrom?, AddrC, Fs1?, Fs2),
    fmrv(Id?, In?, Fs2?).

fmrv(Id, [msg(_, To, punch_to(found(Addr)))|In], Fs) :-
    ground(Id?), Id? =?= To? |
    punch_udp(Addr?),
    fmrv(Id?, In?, Fs?).

fmrv(Id, [msg(_, To, punch_to(not_found))|In], Fs) :-
    ground(Id?), Id? =?= To? |
    fmrv(Id?, In?, Fs?).

fmrv(_, [], _).

procedure lookup_and_send(Constant?, Constant?, AddrResult?, AddrResult,
                          FriendList?, FriendList).
lookup_and_send(Key, Id, OtherAddr, found(Addr?),
    [fe(K, [msg(Id?, K?, punch_to(OtherAddr?))|Out1?], Addr)|Rest],
    [fe(K?, Out1, Addr?)|Rest?]) :-
    Key? =?= K?, ground(Addr?) | true.
lookup_and_send(Key, Id, OtherAddr, KeyAddr?,
    [F|Rest], [F?|Rest1?]) :-
    otherwise |
    lookup_and_send(Key?, Id?, OtherAddr?, KeyAddr, Rest?, Rest1).
lookup_and_send(_, _, _, not_found, [], []).

procedure update_addr(Constant?, Constant?, FriendList?, FriendList).
update_addr(Key, NewAddr, [fe(K, Out?, _)|Fs],
            [fe(K?, Out, NewAddr?)|Fs?]) :-
    Key? =?= K? | true.
update_addr(Key, NewAddr, [F|Fs], [F?|Fs1?]) :-
    otherwise | update_addr(Key?, NewAddr?, Fs?, Fs1).
update_addr(_, _, [], []).
\end{verbatim}

%==============================================================================
\section{IP Cold-Call Invite Redemption}
\label{appendix:ip-cold-call-invites}
%==============================================================================
An IP cold-call peer link is a signed bearer invitation created by agent~$A$ for an as-yet-unknown holder~$B$.  The link is base64url-encoded and is redeemed through a rendezvous server.

\begin{verbatim}
Invite ::= invite(
    A_pk,
    A_nickname,
    RvServer,        %% RV address and public key
    Nonce256,
    ExpiresAt,
    SignatureA
).
\end{verbatim}

\noindent \code{Nonce256} is 256 uniformly random bits.  \code{ExpiresAt} is one hour after creation. \code{RvServer} contains a rendezvous server's public address and key. \code{SignatureA} is an Ed25519 signature by~$A$ over all preceding fields. The nonce is usable exactly once: $A$ records the invite as unused with the rendezvous agent when the link is generated, and the rendezvous agent marks it consumed atomically when validating the redeemer's claim.  Possession of the link authorizes a first-contact attempt; friendship is established only when $A$ accepts the resulting cold call from $B$.

The link is base64-encoded for compact URI transport, allowing the recipient to open it as a deep-link.  The protocol relies on the out-of-band channel to deliver the link to the intended recipient, and relies on the 256-bit nonce, one-hour expiration, and $A$'s signature to prevent guessing, replay, and tampering.

Both $A$ and any holder of the link derive an opaque invite identifier and proof key:
\[
\begin{aligned}
\mathit{InviteId}  &= H(\text{``glp invite id''} \mid \mathit{Nonce256}),\\
\mathit{InviteKey} &= \mathit{HKDF}(\mathit{Nonce256},
    \text{``glp invite proof''} \mid A_{\mathit{pk}} \mid
    \mathit{RvServer} \mid \mathit{ExpiresAt}).
\end{aligned}
\]
The raw nonce is not sent over the IP network.  Network messages use \code{InviteId} plus MACs under \code{InviteKey}, so passive observers and off-path attackers cannot learn or guess the nonce from redemption traffic.

\subsection{Rendezvous-Gated Redemption}
\label{appendix:rv-invite-redemption}

\begin{enumerate}
\item \textbf{Generation and registration.}  $A$ picks a fresh 256-bit secret, computes the invite ID and proof key from it, and assembles the link (public key, nickname, the chosen rendezvous server~$S$'s address and public key, secret, expiration).  $A$ signs the link with its private key and records the invite locally as unused until the expiration.  Over $A$'s existing authenticated connection to~$S$, $A$ registers the invite by sending $A$'s public key, the invite ID, the proof key, and the expiration.  $S$ stores these together with an ``unused'' flag.

\item \textbf{Link delivery.}  $A$ sends the base64url-encoded link to~$B$ out-of-band.

\item \textbf{Cold-call claim.}  $B$ opens the link.  \func{consumePeerLink(uri)} decodes it, verifies $A$'s signature, checks that it is still valid, registers $A$'s public key and nickname locally, and recomputes the invite ID and proof key from the secret.  $B$ opens an encrypted, authenticated connection to~$S$ (using $S$'s public key from the link) and submits a cold-call claim: the invite ID, $B$'s public key, $B$'s nickname, and a short proof computed from the proof key covering $A$'s public key, $B$'s public key, the invite ID, and a fresh nonce of $B$'s own.  $B$ then sends a standard rendezvous \code{available} message naming itself and~$A$.

\item \textbf{Claim validation and forwarding.}  $S$ checks four conditions: the invite ID is registered for~$A$; the expiration is still in the future; the invite is still marked unused; and the proof verifies against the proof key stored in step~1.  On any failure, $S$ rejects the claim.  On success, $S$ atomically flips the invite to ``used'' and forwards $B$'s public key and nickname to~$A$ over $A$'s existing connection to~$S$.

\item \textbf{Acceptance.}  $A$ presents $B$'s public key and nickname to the application layer as a friendship request.  If $A$ accepts the request (which is the friendship decision), $A$ sends~$S$ a standard rendezvous \code{reconnect} message naming~$B$.  In place of the friendship attestation that the reconnect message normally carries, the proof field references the invite ID that $A$ pre-registered with~$S$ in step~1; $S$ matches it against its invite state to authorize the request.

\item \textbf{Signaling.}  $S$ now holds $A$'s \code{reconnect} and $B$'s \code{available} for the same pair of public keys, and the standard rendezvous protocol (Section~\ref{sec:rv-protocol}) takes over with~$S$ as the signaling agent.  $S$ observes both public addresses from the connections it holds and sends each side the other's address.

\item \textbf{UDP path and Noise handshake.}  $A$ and~$B$ each send UDP packets toward the address they received.  Once the UDP path is open, the side designated as initiator (by a deterministic rule based on the public keys) opens a UDX stream, and both sides run the Noise~XX handshake using their Ed25519 keys (Section~\ref{sec:ip-connection}).  The rendezvous server sees signaling metadata and observed addresses, never the encrypted payload.

\item \textbf{Cold-call delivery.}  Only after the Noise-authenticated UDX stream is established does $B$ send the madGLP cold-call payload, which the networking layer delivers to $A$'s index-0 serializer.  The friendship between $A$ and $B$ was established by $A$'s acceptance in step~5.
\end{enumerate}
